The main notation used in the definition of RCSP-SO is summarized in
Table~\ref{tab:notation_RCSP-SO}.
Objective values are evaluated lexicographically by prioritizing total
tardiness and, among solutions with equal total tardiness, minimizing
total sample usage.

\begin{table}[H]
\caption{Sets, parameters and solution measures used in the definition of RCSP-SO}
\label{tab:notation_RCSP-SO}
\centering
\small
\begin{tabular}{cl}
\toprule
\textbf{Notation} & \textbf{Description}\\
\midrule
\textbf{Sets} & \\
\midrule
\(\mathcal{P}\) & Set of products, indexed by \(p\) \\
\(\mathcal{P}^v\) & Subset of products that require nonstandard voltage \\
\(\mathcal{C}\) & Set of chambers, indexed by \(c\) \\
\(\mathcal{C}^h\) & Subset of chambers that have humidity-adjustment capability \\
\(\mathcal{C}^v\) & Subset of chambers that have voltage-adjustment capability \\
\(\mathcal{K}_c\) & Set of stations in chamber \(c\) \\
\(\mathcal{E}\) & Set of environmental (temperature and humidity) configurations, indexed by \(e\) \\
\(\mathcal{E}^H\) & Set of environmental configurations that require 85\% RH \\
\(\mathcal{T}\) & Set of tests, indexed by \(t\) \\
\(\mathcal{T}_p\) & Set of tests required by product \(p\) \\
\midrule
\textbf{Parameters} & \\
\midrule
\(d_p\) & Due date of product \(p\) \\
\(e(t)\) & Environmental configuration required by test \(t\) \\
\(\tau_t\) & Duration of test \(t\) \\
\(K_c\) & Number of stations in chamber \(c\) \\
\(h_c\) & 1 if chamber \(c\) can adjust humidity; 0 otherwise \\
\(v_c\) & 1 if chamber \(c\) can adjust voltage; 0 otherwise \\
\(s_{\min}\) & Lower bound on the sample size of a product \\
\(s_{\max}\) & Upper bound on the sample size of a product \\
\midrule
\textbf{Solution measures} & \\
\midrule
\(C_p\) & Completion time of product \(p\) \\
\(T_p\) & Tardiness of product \(p\), \(T_p=\max\{0,C_p-d_p\}\) \\
\(T^{tot}\) & Total tardiness, \(T^{tot}=\sum_{p\in\mathcal{P}}T_p\) \\
\bottomrule
\end{tabular}
\end{table}
